\documentclass[11pt]{article}
\usepackage[margin=.78in,top=.65in,bottom=.65in]{geometry}
\usepackage{amsmath,amssymb,amsthm,hyperref}
\renewcommand{\normalsize}{\fontsize{10}{12}\selectfont}
\DeclareMathSizes{10}{10}{8}{6}
\hypersetup{colorlinks=true,linkcolor=blue,urlcolor=blue,citecolor=blue}
\setlength{\emergencystretch}{2em}
\newtheorem{theorem}{Theorem}
\newtheorem{lemma}[theorem]{Lemma}
\newcommand{\R}{\mathbb R}
\begin{document}
\begin{center}
{\Large\bfseries A uniform Burnside--Gautschi bound for interpolated gamma values}\\[.5em]
Research proof and Lean formalization\\
October 27, 2026
\end{center}
\begin{abstract}
We prove Conjecture~37 of Schmidlin's paper on Stirling--Gautschi inequalities.
The proposed upper bound holds strictly for every nonnegative base point and,
more generally, every nonnegative interpolation parameter. The argument combines
a classical three-term Stirling remainder bound with a logarithmic entropy estimate.
The remaining rational inequality is certified by a quadratic discriminant whose
coefficients become positive after translation. The Lean proof includes the literal
gamma-function expression and a proof of the classical remainder bound.
\end{abstract}

\section{Statement and reduction}
Following Schmidlin~\cite{S}, put
\[
 \widehat\Pi(x,\alpha)=\Gamma(x+1)(x+1)^\alpha,
 \qquad S_{1/2}(y)=\sqrt{2\pi}\,e^{-1/2}(y+1/2)^{y+1/2}e^{-y},
\]
and $\widehat m(x,\alpha)=\log(\widehat\Pi(x,\alpha)/S_{1/2}(x+\alpha))$.
Conjecture~37 asks for the following weak inequality when $x\ge0$ and
$0\le\alpha\le1$. Corollary~32 of the source proves it for $x\ge1$.
\begin{theorem}\label{main}
For every $x\ge0$ and $\alpha\ge0$,
\[
 \widehat m(x,\alpha)<
 \frac{1}{12(x+\alpha)+18-12e^{1/48}}.
\]
\end{theorem}
All powers and logarithms in this formula have positive bases. No assertion of
optimality is made for the constant.

Set $z=x+1\ge1$, $a=\alpha-1/2\ge-1/2$, and $k=12(e^{1/48}-1)$.
Define the classical Stirling error and an entropy term by
\begin{align}
 R(z)&=\log\Gamma(z)-(z-1/2)\log z+z-\tfrac12\log(2\pi),\label{R}\\
 E(z,a)&=(z+a)\log(1+a/z)-a.\label{E}
\end{align}
Direct algebra gives
\[
 \widehat m(x,\alpha)=R(z)-E(z,a),\qquad
 12(x+\alpha)+18-12e^{1/48}=12(z+a)-k.
\]
The elementary estimates $1+u\le e^u\le(1-u)^{-1}$, $u=1/48$, give
$1/4\le k<1/2$. In particular, the denominator is positive.

\begin{lemma}\label{entropy}
If $z>0$ and $z+a>0$, then $E(z,a)\ge0$. If also $a\ge0$, then
\[
 E(z,a)\ge\frac{a^2}{2z+a}.
\]
\end{lemma}
\begin{proof}
Use $\log t\ge1-t^{-1}$ for the first assertion. For the second, use
$\log(1+t)\ge2t/(2+t)$ with $t=a/z$ and multiply by $z+a$.
\end{proof}

\newpage
\section{An exact positive-polynomial certificate}
The classical Stirling remainder estimate~\cite[\S5.11(ii)]{DLMF} gives
\begin{equation}\label{U}
 R(z)\le U(z):=\frac1{12z}-\frac1{360z^3}+\frac1{1260z^5},\qquad z>0.
\end{equation}
For $z\ge1$, $U(z)\le1/(12z)$. If $-1/2\le a<0$, Lemma~\ref{entropy}
therefore gives immediately
\[
 R(z)-E(z,a)\le\frac1{12z}<\frac1{12(z+a)-k}.
\]
It remains to handle $a\ge0$.

Define
\begin{align*}
 A(z)&=120960z^5,\\
 B(z)&=120960z^6-2520z^5-10080z^4+336z^2-96,\\
 C(z)&=-20160z^5+210z^4+1008z^3-7z^2-288z+2,\\
 F(z)&=420z^5+672z^4-14z^3-192z^2+4z.
\end{align*}
Clearing positive denominators yields the identity
\begin{equation}\label{cert}
 \frac1{12(z+a)-1/4}-U(z)+\frac{a^2}{2z+a}
 =\frac{A(z)a^3+B(z)a^2+C(z)a+F(z)}
 {2520z^5(a+2z)(48a+48z-1)}.
\end{equation}
Both $B(z)$ and $\Delta(z):=4B(z)F(z)-C(z)^2$ are strictly positive for
$z\ge1$. Indeed, writing $z=1+t$, $t\ge0$, their coefficients in descending
order are respectively
\[
 120960,\ 723240,\ 1791720,\ 2353680,\ 1729056,\ 673512,\ 108600
\]
and
\begin{gather*}
 203212800,\ 2149822080,\ 10299502080,\ 29404384380,\\
 55375719840,\ 71883106764,\ 65187512040,\ 40869159443,\\
 17049349100,\ 4345739674,\ 553547260,\ 16630775.
\end{gather*}
Thus
\[
 4B(z)\bigl(B(z)a^2+C(z)a+F(z)\bigr)
 =(2B(z)a+C(z))^2+\Delta(z)>0.
\]
The cubic term $A(z)a^3$ is nonnegative. Formula~\eqref{cert} and
Lemma~\ref{entropy} now imply
\[
 R(z)-E(z,a)\le U(z)-\frac{a^2}{2z+a}
 <\frac1{12(z+a)-1/4}\le\frac1{12(z+a)-k},
\]
which proves Theorem~\ref{main}.

The polynomial identities above are checked in Lean by ring normalization.
Their coefficient positivity, the quadratic identity, and every denominator
sign are also proved in Lean. The integer lists in this paper are not trusted
inputs to an external numerical computation.

\newpage
\section{The classical remainder bound in the formal proof}
For completeness, the certificate proves~\eqref{U} from elementary logarithmic
bounds and the pinned library's integer Stirling theorem. The gamma recurrence gives
\[
 R(z)-R(z+1)=h(z):=(z+1/2)\log(1+1/z)-1.
\]
For $v=(2z+1)^{-1}\in(0,1)$, the atanh-series remainder estimate gives
\[
 h(z)\le H(z):=\frac{v^2}{3}+\frac{v^4}{5}+\frac{v^6}{7}
                    +\frac{v^8}{1-v^2}.
\]
A second exact identity is
\[
 U(z)-U(z+1)-H(z)=
 \frac{112z^8+448z^7+1295z^6+2317z^5+2200z^4+1061z^3+253z^2+34z+2}
 {2520z^5(z+1)^5(2z+1)^6}>0.
\]
Consequently $R(z)-U(z)$ increases under $z\mapsto z+1$. Its limit is zero,
so $R(z)\le U(z)$.

Here the real-variable limit is itself proved, rather than assumed as a
certificate hypothesis. Log-gamma convexity gives, for integers $n\ge2$ and
$0\le r\le1$,
\[
 \log\Gamma(n)+r\log(n-1)\le\log\Gamma(n+r)
                       \le\log\Gamma(n)+r\log n.
\]
Using $u/(1+u)\le\log(1+u)\le u$ yields
\[
 R(n)-\frac2{n-1}\le R(n+r)\le R(n)+\frac1n.
\]
At integers, $R(n)$ is the logarithm of Mathlib's Stirling sequence minus
$\tfrac12\log\pi$, and therefore tends to zero. Taking $n=\lfloor z\rfloor$
proves the required limit for real $z\to\infty$.

\section{Formal scope and prior work}
The Lean endpoint \texttt{Schmidlin.source\_conjecture\_372858} is the complete
source conjecture. The stronger endpoint \texttt{source\_strict\_bound} quantifies
over every $x\ge0$ and $\alpha\ge0$. The definitions use the actual gamma function,
real powers, and the literal logarithm of the interpolated-value/Burnside ratio.
The logarithmic expansion and all analytic and rational bridges are proved.
No classical Stirling or entropy fact is claimed as a new discovery.

The source's established $x\ge1$ range is credited to Schmidlin. A bounded
preliminary search found no existing resolution of the remaining conjecture;
it does not certify worldwide priority.

\begin{thebibliography}{9}
\bibitem{S} M.~Schmidlin, \emph{On Gautschi \& Stirling identities, asymptotics
and inequalities for the Pi (or Gamma) function}, arXiv:2601.21906v1 (2026),
Conjecture~37, pp.~28--29. \url{https://arxiv.org/abs/2601.21906v1}.
\bibitem{DLMF} NIST Digital Library of Mathematical Functions,
\S5.11(i)--(ii), Stirling's series and remainder bounds.
\url{https://dlmf.nist.gov/5.11}.
\bibitem{M} The Mathlib Community, \emph{The Lean mathematical library},
pinned revision \texttt{d13f23b723b8a846827a245b89c10fc7d3f11612};
\texttt{Stirling.lean}, \texttt{Gamma/BohrMollerup.lean}.
\end{thebibliography}
\end{document}
